\documentclass[11pt,a4paper]{article}
\usepackage[margin=24mm,headheight=14pt,includeheadfoot]{geometry}
\usepackage[T1]{fontenc}
\usepackage{amsmath,amsthm}
\usepackage{newtxtext,newtxmath,microtype}
\usepackage{booktabs,longtable,array,graphicx,pdflscape,caption,fancyhdr,xcolor,hyperref,placeins,float}
\definecolor{ink}{HTML}{23384D}
\hypersetup{colorlinks=true,linkcolor=ink,urlcolor=ink,citecolor=ink,
 pdftitle={Variance learning rates and backbones in KSIVI},pdfauthor={StatComp research investigation}}
\newtheorem{proposition}{Proposition}[section]
\newtheorem{corollary}[proposition]{Corollary}
\newcommand{\E}{\mathbb E}
\newcommand{\Var}{\operatorname{Var}}
\newcommand{\Cov}{\operatorname{Cov}}
\newcommand{\tr}{\operatorname{tr}}
\newcommand{\diag}{\operatorname{diag}}
\setlength{\parskip}{4pt}\setlength{\parindent}{1em}
\setlength{\tabcolsep}{4pt}\renewcommand{\arraystretch}{1.16}
\setlength{\emergencystretch}{1.5em}
\captionsetup{font=small,labelfont=bf,labelsep=period}
\pagestyle{fancy}\fancyhf{}
\fancyhead[L]{\small\itshape KSIVI: variance learning rates}
\fancyhead[R]{\small Atomic comparisons and mechanisms}
\fancyfoot[C]{\thepage}\renewcommand{\headrulewidth}{0.3pt}
\title{Can a Slower Variance Network Stabilize KSIVI?\\
\large Learning rates, shared features, and optimization budgets on the X-shaped target}
\author{StatComp research investigation}\date{12 October 2026}
\begin{document}\maketitle
\begin{abstract}
@ABSTRACT@
\end{abstract}

\section{Question, findings, and the scope of the evidence}
The motivating suggestion is to reduce the variance learning rate while retaining
the mean learning rate. That suggestion has two distinct implementations:
slowing only a variance output layer on shared features, and slowing an entire
independent variance network. They need not have the same effect. We return to
the original StatComp model and training policy; the authors' alternative
defaults are not substituted for that policy in this campaign.

@FINDINGS@

\begin{center}\small
@SUMMARY@
\end{center}

The statistical unit is an independently trained model. Monte Carlo samples
from one checkpoint do not count as additional training replications. Discovery
seeds select hyperparameters; held-out seeds assess those choices. Descriptive
means and sample standard deviations are reported, together with paired
per-seed outcomes. This is an exploratory investigation, not a preregistered
confirmatory significance test. No statement below establishes a universal
ordering of optimizers, architectures, or gradient estimators.

\section{The common model and loss}
Let $Z\sim N(0,I_2)$ and $U\sim N(0,I_2)$ be independent. StatComp samples
\[
 X=\mu_\phi(Z)+\sqrt{v_\phi(Z)}\odot U,\qquad
 v_\phi(Z)=\max\{10^{-4},\operatorname{softplus}(a_\phi(Z))\}.
\]
The mean uses two width-128 SiLU hidden layers. The global model replaces
$v_\phi(Z)$ by two learnable variance scalars. At initialization all matched
comparisons have the same mean function and constant variance
$\operatorname{softplus}(1)=1.3132616$, except explicitly labeled initialization
ablations. The effective mixing dimension is two: the production target binding
overrides the raw target-configuration entry.

Write $s_p=\nabla\log p$ and $s_c(X,Z)=\nabla_X\log q_\phi(X\mid Z)$.
The conditional Gaussian layer returns its negative score
$-s_c=U/\sqrt v$. With two independent batches of size $N=128$, the original
loss estimator is
\begin{equation}
 \widehat L_c=\frac1{N^2}\sum_{i,j} k_h(X_i,Y_j)
 \left(\alpha_t s_p(X_i)+\frac{U_i}{\sqrt{v(Z_i)}}\right)^T
 \left(\alpha_t s_p(Y_j)+\frac{U'_j}{\sqrt{v(Z'_j)}}\right).
 \label{eq:original}
\end{equation}
For a kernel fixed independently of these samples, the Fisher identity
$\E[s_c\mid X]=s_q(X)$ shows that its expectation is the squared KSD between
$q_\phi$ and $p_{\alpha_t}\propto p^{\alpha_t}$, under the usual integration
and differentiation conditions. This is the tractable conditional-score
formulation underlying KSIVI \cite{ksivi}. A larger conditional family need
not have a larger minimum population loss; our question concerns finite
stochastic optimization.

For the Gaussian kernel $k_h(x,y)=\exp[-\|x-y\|^2/(2h^2)]$, integration by
parts gives the conventional Stein form
\begin{equation}
 \widehat L_S=\frac1{N^2}\sum_{i,j} k_h(X_i,Y_j)
 \left[s_i^Ts'_j+\frac{(s_i-s'_j)^T(X_i-Y_j)}{h^2}
 +\frac2{h^2}-\frac{\|X_i-Y_j\|^2}{h^4}\right],
 \label{eq:stein}
\end{equation}
where $s_i=\alpha_t s_p(X_i)$. For independently fixed $h$, the two forms
have the same population objective, but different finite-batch gradients.
Neither form has a universally smaller gradient covariance.

Production training fits
$h^2=\operatorname{median}_{ij}\|X_i-Y_j\|^2/[2\log(N+1)]$ on each batch.
Original-loss runs differentiate this fitted width. Stein runs hold it fixed
during backpropagation, matching the earlier investigation. Thus the default
Stein comparison changes both estimator and bandwidth differentiation. We
include original-loss bandwidth-detachment controls and independently fixed
bandwidth probes; the fixed-kernel population identity is not asserted for a
sample-adaptive differentiated statistic.

\subsection{Relation to the original KSIVI guarantees}
The KSIVI analysis assumes controlled model derivatives and a uniform positive
lower bound on component standard deviations (Assumption 4.3), then establishes
gradient-variance bounds and an SGD stationarity guarantee under its full
assumptions (Theorems 4.7--4.8) \cite{ksivi}. Our small-variance limit leaves a
fixed uniformly nondegenerate parameter domain; it does not contradict those
results. The implemented $10^{-4}$ variance floor truncates that limit but can
still permit large noise constants. Stationarity also does not imply a small
transport error. Adam, annealing, and fitted bandwidths require further care
before applying a fixed-objective SGD theorem to these runs.

There is an additional exact boundary on this particular target: its log-density
does not have a globally bounded Hessian. In fact, with $c=1.8/0.76$,
\[
 \log p(x,y)=C-\frac{x^2+y^2}{0.76}+\log\cosh(cxy),\qquad
 \partial_{yy}\log p(R,0)=-\frac2{0.76}+c^2R^2.
\]
The last expression is unbounded as $R\to\infty$, so the target does not satisfy
the global Hessian condition in Assumption 4.2. Spatial autograd of the actual
target implementation checks the formula. This is a limitation on a direct
theorem application, not an explanation of the ordering of methods: all
methods in this campaign share the same target.

\section{What a variance learning rate can and cannot control}
\subsection{A head-only learning rate leaves another path active}
For a shared backbone $h_b(Z)$, write
\[
 \mu(Z)=M_\psi h_b(Z)+m_\psi,\qquad
 a(Z)=W_\omega h_b(Z)+c_\omega.
\]
The backbone gradient decomposes exactly as
\begin{equation}
 g_b=J_b\mu^T\partial_\mu\widehat L+J_b a^T\partial_a\widehat L
       =g_b^{\mu}+g_b^a.
 \label{eq:paths}
\end{equation}
A smaller learning rate for $(W,c)$ leaves the second term in the backbone's
gradient active at the backbone learning rate. Mean-driven backbone updates
also change the features read by the variance head.

\begin{proposition}[Exact functional drift decomposition]
For a fixed probe input, let $h'=h+\Delta h$, $W'=W+\Delta W$, and
$c'=c+\Delta c$ denote the values before and after one optimizer step. Then
\begin{equation}
 \Delta a=\underbrace{\Delta W h+\Delta c}_{\text{head}}
          +\underbrace{W\Delta h}_{\text{backbone}}
          +\underbrace{\Delta W\Delta h}_{\text{interaction}}.
 \label{eq:drift}
\end{equation}
Consequently, freezing a trained variance head does not generally freeze its
predicted variance on a changing shared backbone.
\end{proposition}
\begin{proof}
Expand $(W+\Delta W)(h+\Delta h)+(c+\Delta c)-(Wh+c)$.
When $\Delta W=\Delta c=0$, the remaining term is $W\Delta h$, which need
not vanish. At our constant initialization $W=0$, this route is initially
absent; after the head has learned, it can be present.
\end{proof}

At a fixed state, a first-order SGD calculation gives
$\Delta a\simeq-\eta_v J_\omega a\,g_\omega-\eta_bJ_b a\,g_b$.
As $\eta_v\to0$ with $\eta_b$ fixed, its conditional covariance tends to
$\eta_b^2J_b a\,\Cov(g_b)J_b a^T$ when this matrix is nonzero. Cross
covariances matter at finite $\eta_v$; the formula is not a universal positive
lower bound for every model or state. It does show why head LR alone cannot
guarantee control of functional variance noise.

For an independent variance network $a_\omega(Z)$, all of its parameters
receive $\eta_v$. To first order, both functional drift and stochastic
standard deviation scale with $\eta_v$ at a fixed state. This is genuine
control of the whole variance function, while also slowing its deterministic
learning. The gradient-path and trained-head-freeze probes test these routes
directly rather than inferring them from the final sample plots.

Our stop-gradient ablation detaches the shared features only along the
variance path. It removes $g_b^a$ but does not prevent $h_b$ from moving under
$g_b^\mu$. Moreover, this deliberately omits part of the shared model's
gradient: it is a routing heuristic, not an unbiased gradient of the original
shared objective. Separating backbones retains the full gradient of the new
parameterization.

\subsection{Different rates do not create gradient-estimation bias}
\begin{proposition}[Block learning rates and descent]
Let $L$ have a $\beta$-Lipschitz gradient, and let a stochastic gradient
$\widehat g$ at a fixed state have mean $g=\nabla L$ and covariance $\Sigma$.
For a deterministic positive block-diagonal learning-rate matrix $D$,
\[
 \E[D\widehat g]=Dg,
\]
and
\begin{equation}
 \E[L(\theta-D\widehat g)]\le L(\theta)-g^TDg
 +\frac\beta2\{g^TD^2g+\tr(D^2\Sigma)\}.
 \label{eq:descent}
\end{equation}
\end{proposition}
\begin{proof}
Apply the smoothness inequality to $-D\widehat g$, take expectations, and use
$\E[\widehat g\widehat g^T]=gg^T+\Sigma$.
\end{proof}

Thus changing a block LR preconditions the update; it does not by itself bias
the underlying gradient estimator. For that block, the expected first-order
movement scales as $\eta_v$, while the covariance of the SGD update scales as
$\eta_v^2$. Smaller steps can suppress noisy motion but also delay progress.
This is a local SGD statement. Adam's moment estimates are random and
history-dependent, so it is not claimed as an Adam convergence theorem.

For Adam, the learning rate multiplies the normalized direction
$\widehat m/(\sqrt{\widehat s}+\varepsilon)$ \cite{adam}. At fixed parameters,
historical moments, and a fresh stochastic gradient, scaling that block's LR
scales its parameter update. Multiplying its gradient by a constant is not
equivalent: the first and second moments approximately cancel the scale when
$\varepsilon$ is negligible. We therefore implement actual parameter groups,
not gradient multipliers. The legacy combined output layer is split into two
parameter blocks and concatenated for the same forward Linear operation.

\subsection{Separating backbones does not imply half the convergence speed}
Both networks receive gradients from the same loss and both are updated at
every iteration. There is no factor $1/2$ in either block's gradient. A simple
counterexample rules out a necessary twofold slowdown: for
$L(m,a)=(m^2+a^2)/2$, joint SGD and two simultaneously updated separate
blocks both obey $(m_{t+1},a_{t+1})=(1-\eta)(m_t,a_t)$ exactly. A separate
network can require more computation and can change conditioning and capacity;
those facts do not imply a fixed change in the number of optimizer steps.

\section{Why an arbitrarily slow variance network is not a complete solution}
\subsection{The initial scale must be able to shrink}
The X-shaped target is
\[
 p=\tfrac12N(0,\Sigma_+)+\tfrac12N(0,\Sigma_-),\qquad
 \Sigma_\pm=\begin{pmatrix}2&\pm1.8\\\pm1.8&2\end{pmatrix}.
\]
Its component eigenvalues are $3.8$ and $0.2$. If $S$ is an equiprobable sign,
$T\sim N(0,1)$, and $U\sim N(0,I_2)$ independently, then
\[
 X=\sqrt{1.8}\,T(1,S)^T+\sqrt{0.2}\,U
\]
has exactly this distribution. A global component variance of $0.2$ is
therefore sufficient with an unrestricted mixing distribution of means; a
finite continuous neural network approximates the discrete branch selection.

\begin{proposition}[An obstruction for frozen large component variance]
No mixing distribution of means convolved with $N(0,vI_2)$ can equal this
target when $v>0.2$.
\end{proposition}
\begin{proof}
In direction $w=(1,-1)^T/\sqrt2$, the target characteristic function is
$\varphi_p(tw)=\tfrac12e^{-0.1t^2}+\tfrac12e^{-1.9t^2}$. A representation
$X=M+\sqrt vU$ would require
$\varphi_M(tw)=\varphi_p(tw)e^{vt^2/2}$, which is unbounded as
$|t|\to\infty$ if $v>0.2$. Every characteristic function has modulus at most
one, a contradiction.
\end{proof}

More quantitatively, any such frozen-variance model satisfies
\[
 W_2(w^Tq,w^Tp)\ge W_1(w^Tq,w^Tp)
 \ge\sup_{t>0}\frac{[\varphi_p(tw)-e^{-vt^2/2}]_+}{t}.
\]
The last inequality follows from the Lipschitz constant $t$ of
$x\mapsto e^{itx}$ and $|\varphi_q(tw)|\le e^{-vt^2/2}$.
For the initial variance 1.31326, our numerical one-dimensional bound is
about 0.13535. It is a bound on this projection, not the averaged sliced
metric, and does not explain the full magnitude of an escaping training run.

This distinguishes a hard floor of $0.2$ from freezing the initial variance
near $1.3$. The floor permits shrinkage to the target's transverse scale;
freezing the initial value prevents that adaptation. Initialization at $0.2$
and its frozen control test this explanation directly.

\subsection{A finite-budget obstruction for a very slow linear variance head}
There is also a deterministic Adam bound that can make under-adaptation precise.
Consider Adam without weight decay, initialized with zero moments, with
$\beta_1^2<\beta_2$ and nonnegative numerical $\varepsilon$. For any coordinate,
\begin{equation}
 \left|\frac{\widehat m_t}{\sqrt{\widehat s_t}+\varepsilon}\right|
 \le C:=\frac{1-\beta_1}{\sqrt{(1-\beta_2)(1-\beta_1^2/\beta_2)}}.
 \label{eq:adam-bound}
\end{equation}
For our betas, $C=7.2703$. This bounds the coordinate update by $C\eta_t$
for every finite gradient sequence; no stochastic independence is assumed.

\begin{proposition}[Reachability bound for a linear variance model]
Let $a_j(Z)=w_j^TZ+b_j$, initially $w_j=0$ and $b_j=1$, with every variance
parameter assigned LR $r\eta_t$. Put $\delta=C r\sum_{t<T}\eta_t$.
On the latent region $B_M=\{\|Z\|_1\le M\}$, each variance is at least
$\ell=\operatorname{softplus}(1-(1+M)\delta)$. Since
$P(B_M)=\operatorname{erf}(M/2)^2=:P_M$, regardless of the trained mean,
\begin{equation}
 W_2(w^Tq_T,w^Tp)\ge
 \sup_{t>0}\frac{[\varphi_p(tw)-P_Me^{-\ell t^2/2}-(1-P_M)]_+}{t},
 \label{eq:reachable}
\end{equation}
where $w=(1,-1)^T/\sqrt2$. A positive right-hand side rules out an exact fit
within the specified optimization budget.
\end{proposition}
\begin{proof}
For completeness, Cauchy--Schwarz applied to the two weighted gradient moments
gives the squared bound
\[
 C_t^2=\frac{(1-\beta_1)^2}{(1-\beta_1^t)^2}
 \frac{1-\beta_2^t}{1-\beta_2}
 \frac{1-(\beta_1^2/\beta_2)^t}{1-\beta_1^2/\beta_2}.
\]
The inequality $(1-x)(1-y)\le(1-\sqrt{xy})^2$, with
$x=\beta_2^t$ and $y=(\beta_1^2/\beta_2)^t$, shows $C_t\le C$.
Summing coordinate updates gives $|w_{jk}|\le\delta$ and $|b_j-1|\le\delta$,
so $a_j(Z)\ge1-(1+M)\delta$ on $B_M$. Orthogonal rotation of a standard
bivariate Gaussian gives $P(B_M)=\operatorname{erf}(M/2)^2$.
Conditional noise in direction $w$ has variance $(v_1+v_2)/2\ge\ell$ on
$B_M$. Therefore
$|\varphi_{q_T}(tw)|\le P_Me^{-\ell t^2/2}+(1-P_M)$, independently of
the conditional means. The Lipschitz characteristic-function argument from
the previous subsection gives \eqref{eq:reachable}.
\end{proof}

For $r=0.001$, $T=50K$, and $M=3$, the coordinate displacement is at most
about $0.0723$ and $P_M\simeq0.9334$. The resulting projection bound is
strictly positive. Retaining the original decay through 200K scarcely changes
it. The recorded numerical checks give the full bounds for the restart policy
and additional rates. These are bounds for the linear variance architecture,
in exact arithmetic; they are not bounds for an unrestricted changing shared
feature map. Linear-head runs at ratios 0.001 and 0.003 are included to test
this predicted finite-budget limitation. The result shows why a universally
beneficial prescription to decrease the variance LR cannot hold, even before
considering stochastic gradient variance.

\subsection{Representation noise can grow without a worse marginal}
The conditional score has the exact coordinate second moment
$\E[s_{c,j}^2]=\E[1/v_j(Z)]$. This is a diagnostic of representation noise,
not a bound on every parameter-gradient variance or a proof of poor marginal
fit. The following solvable example makes that distinction explicit.

\begin{proposition}[An exact finite-batch gradient-noise example]
Let $p=N(0,1)$ and, for $0<v<1$, represent $q_v=p$ as
$X=\sqrt{1-v}\,E+\sqrt v\,U$, where $E,U$ are independent standard Normals.
Take the constant kernel $k=1$, two independent batches of size $N$, and
differentiate with respect to $r=\log v$ while preserving this representation.
Both population objectives are zero for every $v$, but their stochastic
gradient variances are
\begin{align}
 \Var(\widehat g_c)&=\frac{2+v-v^2}{2N^2v^2},\\
 \Var(\widehat g_S)&=\frac{v}{2(1-v)N^2}.
 \label{eq:proxy}
\end{align}
\end{proposition}
\begin{proof}
Put $F=-X+U/\sqrt v$ and $B=\partial_rF$. Direct differentiation gives
\[
 F=-\sqrt{1-v}E+\frac{1-v}{\sqrt v}U,\quad
 B=\frac{v}{2\sqrt{1-v}}E-\frac{1+v}{2\sqrt v}U.
\]
Hence $A=\Var(F)=(1-v)/v$,
$C=\Var(B)=(1+v-v^2)/[4v(1-v)]$, and
$D=\Cov(F,B)=-1/(2v)$. The original gradient is
$\bar B\bar F'+\bar F\bar B'$, with independent primed batches, so its
variance is $2(AC+D^2)/N^2$, giving the first expression.
For the Stein form, $k=1$ removes all spatial kernel derivatives and the loss
is $\bar X\bar X'$. Here $\Var(X)=1$,
$\Var(\partial_rX)=v/[4(1-v)]$, and
$\Cov(X,\partial_rX)=0$. The same calculation gives the second expression.
\end{proof}

The constant kernel is an analytically solvable proxy, not the characteristic
Gaussian kernel used in training. Independent fixed-bandwidth RBF simulations
with the same unchanged marginal check whether the mechanism also appears
with the training kernel family. The formulas prove a possible noise mechanism,
not a universal RBF variance ordering: near $v=0.8$ even the constant-kernel
ordering reverses. A small learning rate changes update noise, while Stein
integration removes the explicit component-score terms themselves.

\begin{proposition}[Small-variance asymptotics for the Gaussian kernel]
Let $p=q_v=N(0,I_d)$ be represented as
$X_v=\sqrt{1-v}\,E+\sqrt v\,U$ with independent standard Gaussian vectors,
and use two independent batches of fixed size $N$. For fixed $h>0$ and
$r=\log v$, the original Gaussian-kernel gradient satisfies
\begin{equation}
 \lim_{v\downarrow0}v^2\Var(\partial_r\widehat L_c)
 =\frac{d}{N^2}\left(1+\frac4{h^2}\right)^{-d/2}>0,
 \qquad \Var(\partial_r\widehat L_S)=O(v).
 \label{eq:rbf-asymptotic}
\end{equation}
Both population gradients are zero for every $v$. Thus the noise divergence
also occurs with a characteristic Gaussian kernel, without changing the
marginal distribution or the population objective.
\end{proposition}
\begin{proof}
For the original estimator $F_v=-X_v+U/\sqrt v$,
$F_v=U/\sqrt v+O_{L^p}(1)$ and
$\partial_r F_v=-U/(2\sqrt v)+O_{L^p}(\sqrt v)$ for every finite $p$.
Also $\partial_rX_v=O_{L^p}(\sqrt v)$. Gaussian kernel derivatives are
bounded, so differentiation of its factor contributes only
$O_{L^2}(v^{-1/2})$ to the gradient. Therefore
\[
 v\partial_r\widehat L_c\ \longrightarrow\
 -\frac1{N^2}\sum_{i,j}k_h(E_i,E'_j)U_i^TU'_j
 \quad\text{in }L^2.
\]
All cross terms except identical $(i,j)$ vanish after conditioning on the
latent $E$'s and taking expectations over independent centered $U$'s.
The limiting variance is
$d\E[k_h(E,E')^2]/N^2$.
Since $E-E'\sim N(0,2I_d)$, this kernel moment equals
$(1+4/h^2)^{-d/2}$. The Stein integrand and its first spatial derivatives
have uniformly bounded Gaussian moments, and its only $r$ dependence is
through $X_v,Y_v$. Its gradient is consequently $O_{L^2}(\sqrt v)$,
which gives the stated variance order. The $L^2$ limit also justifies passage
of the mean and variance; the common population loss is identically zero.
\end{proof}

For $d=2$, $N=128$, and $h=0.75$, the leading coefficient is approximately
$1.5050\times10^{-5}$. Numerical simulations compare $v^2$ times the original
gradient variance with this coefficient, rather than relying only on a visual
log--log slope. This theorem concerns a specified marginal-preserving direction;
it is not a lower bound for every parameter direction in a trained neural model.

\section{Experimental protocol and atomic controls}
All primary results use 50,000 updates. The fixed policy is Adam
$(0.9,0.999)$, mean LR $10^{-3}$, two batches of 128, StepLR multiplier 0.9
per 1,000 updates, and score annealing
$\alpha_t=0.1+0.9\min(t/25000,1)$. The Gaussian distance arithmetic remains
the StatComp implementation. Diagnostic RNG is isolated. We omit a
diagnostic-only bandwidth-to-CPU synchronization in the fast update; equality
of the loss, gradients, and production optimizer update is checked on CPU and
GPU. No gradient clipping, EMA, or log-density regularization is introduced.

Discovery sweeps variance/mean LR ratios
$\{0,0.001,0.003,0.01,0.03,0.1,0.3,1,3\}$ with seeds 42--44, for shared
and separate width-128 two-layer backbones. For the shared model, only the
variance output parameters receive this rate; for a separate model, the
entire variance network receives it. Both networks are updated every step.
Separate width-32, one-hidden-layer, linear, and shared stop-gradient
alternatives are each tested at ratios $\{0.01,0.1,1\}$. Initial mean
functions, constant variances, and training-noise seeds are paired.

The minimum mean discovery SW$_2$ at a positive LR selects one shared, one
separate, and one alternative configuration. That selection is recorded
before held-out seeds 45--49. Held-out controls include equal-LR shared and
separate models, global variance, Stein loss, and a $0.2$ floor. Further
paired ablations change only annealing, initial variance, bandwidth detachment,
or estimator/floor for the selected head LR. Every such change is labeled.

Training diagnostics use 10,000 samples and 128 projection directions.
Independent final evaluation uses 20,000 samples and 128 directions, plus a
second 10,000-sample check. Forward KL uses 2,048 exact target points and
an empirical Gaussian mixture for $q$: mixture sizes 8,192, 16,384, and
65,536 are checked, and 262,144 is used when the last change exceeds 0.02.
The last change is retained; the pointwise Monte Carlo standard error does
not include all mixture-approximation uncertainty. Held-out KSD uses 16
independent cross-batch estimates, 256 samples per batch, at independently
fixed bandwidths $0.5,0.75,1.5$. Negative finite-sample estimates are retained.

\subsection{Why update counts and learning-rate budgets are both reported}
With $\eta_t=10^{-3}0.9^{\lfloor t/1000\rfloor}$, the first 50,000 updates
have total nominal LR
\[
 S_{50K}=1000\cdot10^{-3}\frac{1-0.9^{50}}{1-0.9}=9.94846.
\]
The sum through 200,000 is almost 10, adding only 0.518\% to $S_{50K}$.
These sums are not actual Adam path lengths, but they show why extra decayed
updates need not compensate for a smaller variance LR. We fork exact 50K
checkpoints, including Adam and RNG state, to two labeled continuations:
retained StepLR decay, or a restart at mean LR $10^{-4}$ with the original
variance/mean LR ratio, held constant. The latter adds nominal mean LR 15
over the next 150K updates. Its improvement cannot be attributed to update
count alone. All continuations are compared at 100K and 200K, and no run
exceeds 200K updates.

\section{Results and ablations}
@RESULTS@

\section{Mechanistic evidence and boundaries}
@MECHANISMS@

\section{Claim-to-evidence audit and practical interpretation}
@CLAIMS@

\subsection{Limitations}
Discovery minima are selected estimates, so held-out results carry the main
weight in recommendations. Five held-out seeds give limited precision for
rare training failures. The investigation covers one two-dimensional target,
one batch size, one mean architecture, one variance transform, and finite
training budgets. Shared versus separate backbones change parameterization
and capacity; narrow, shallow, and linear variants examine but do not eliminate
all such differences. Gradient-path correlations and inverse-variance moments
are state diagnostics, not alone proofs of causal responsibility. The freeze
intervention, paired training ablations, and exact formulas provide different
forms of evidence and are identified separately. Adaptive bandwidths and
annealing limit direct application of fixed-population SGD theory; no global
convergence theorem for the implemented Adam procedure is claimed.

\appendix
\section{Complete numerical comparisons}
@APPENDIX_TABLES@

\clearpage
\begin{landscape}
\section{Final sample distributions}
@SAMPLE_FIGURES@
\end{landscape}

\section{Reproduction, validation, and artifact inventory}
@REPRODUCIBILITY@

\begin{thebibliography}{9}
\bibitem{ksivi} Z. Cheng, L. Yu, T. Xie, S. Zhang, and C. Zhang.
\emph{Kernel Semi-Implicit Variational Inference}. ICML, 2024.
\url{https://arxiv.org/abs/2405.18997}.
\bibitem{adam} D. P. Kingma and J. Ba.
\emph{Adam: A Method for Stochastic Optimization}. ICLR, 2015.
\url{https://arxiv.org/abs/1412.6980}.
\end{thebibliography}
\end{document}
